\documentclass[10pt,a4paper]{amsart}
\usepackage[margin=19mm]{geometry}
\usepackage{amsmath,amssymb,mathtools}
\usepackage[scr=rsfs,cal=euler]{mathalfa}
\usepackage{xfrac}
\usepackage[unicode,hidelinks,bookmarks=false]{hyperref}
\newcommand{\ie}{i.e.\ }
\newcommand{\cf}{cf.\ }
\newcommand{\resp}{respectively\ }
\newcommand{\Subsection}{\subsection}
\providecommand{\nobreakdash}{\nobreak\hbox{-}}
\setlength{\emergencystretch}{2em}
\multlinegap0pt
\usepackage{algorithm}\usepackage{algpseudocode}
\usepackage{amssymb}
\usepackage{xcolor}
\newtheorem{assumption}{Assumption}
\newtheorem{theorem}{Theorem}
\newtheorem{lemma}{Lemma}
\usepackage{cleveref}
\crefname{assumption}{Assumption}{Assumptions}
\crefname{theorem}{Theorem}{Theorems}
\crefname{lemma}{Lemma}{Lemmas}
\newcommand{\N}{\mathbb{N}}
\newcommand{\R}{\mathbb{R}}
\DeclareMathOperator{\cl}{cl}
\newcommand{\norm}[1]{\left\|#1\right\|}
\DeclareMathOperator{\ri}{ri}
\title{Achieving Directional-Stationarity from a Single Random Direction Step}
\author{Dan Greenstein, Nadav Hallak}
\date{Source: arXiv:2605.22045v1 (CC BY 4.0)}
\begin{document}
\maketitle

\noindent\emph{Source and licence.} Achieving Directional-Stationarity from a Single Random Direction Step. Version arXiv:2605.22045v1 (CC BY 4.0), distributed by arXiv under the Creative Commons Attribution 4.0 International licence, http://creativecommons.org/licenses/by/4.0/, as recorded in the arXiv licence metadata for this exact version and shown on the abstract page https://arxiv.org/abs/2605.22045v1. Full licence notice: \texttt{licenses/arxiv-2605.22045-CC-BY-4.0.txt}.\par\smallskip

\noindent\textit{Editorial scope.} Counted unit: the article's theorem thm:random (source0.tex lines 323--329). The article's complete original proof of it is delivered with it (source0.tex lines 554--626, one \texttt{proof} environment of 73 lines, which the article places away from the statement). No mathematics is written, repaired or re-derived: the statement and the proof are the article's own lines, in the article's own notation, and no retained line is substituted or re-flowed.  Retained uncounted as the source-local context the proof needs: lines 104--106; lines 107--109; lines 294--304; lines 361--366; lines 471--477; lines 524--526.  Nothing was omitted from any retained span, so the package carries no omission record.  No retained line contains a citation, so the wrapper carries no bibliography and no bibliography entry.  No retained line contains a figure environment, an image-inclusion command or drawing code, so no image is delivered and the package declares no asset dependency.  Editorial formatting only, in addition: an AMS \texttt{amsart} preamble that restates the article's own declarations and macros used by the retained lines, namely the environments \texttt{assumption}, \texttt{theorem}, \texttt{lemma} on separate counters, together with the standard packages those lines need.  The retained environments print in this document as Assumption 1, Theorem 1, Lemma 1 in order of appearance, whereas the article prints the same items with the numbering of its own full document.  The complete native source of the article as distributed in the arXiv e-print for this exact version is bundled unchanged as \texttt{sources/201105/source0.tex}.
\begin{assumption}[A1]\label{ass:A1}
$h$ is locally Lipschitz and bounded below on $C$.
\end{assumption}

\begin{assumption}[A2]\label{ass:A2}
The directional derivative $h'(x;\, d) = \lim_{t \downarrow 0}(h(x + td) - h(x))/t$ exists for every $x \in C$ and $d \in \R^n$.
\end{assumption}

\begin{algorithm}
\caption{Random exploration procedure implementation}
\label{alg:random}
\begin{algorithmic}[1]
\Require $x^0 \in C$; oracle $\mathcal{O}_C$ over $C$; $\gamma > 0$, $r \in (0,\infty)$.
\For{$k = 0, 1, 2, \ldots$}
  \State Set $y^{k+1} \leftarrow \mathrm{REP}(x^k)$ and  $z^{k+1} \leftarrow \mathcal{O}_C(x^k)$
  \State $x^{k+1} \leftarrow \arg\min_{u\,\in\,\{y^{k+1},\,z^{k+1}\}} h(u)$
\EndFor
\end{algorithmic}
\end{algorithm}

\begin{lemma}[Joint density of exploration direction, proximity, and step size]\label{lem:BC-directions}
There exists an event $\Omega_{BC}$ with $\mathbb{P}(\Omega_{BC})=1$ such that, on $\Omega_{BC}$, for every $0 < a < b \leq r$, every $\rho > 0$, every accumulation point $x^*$ of $\{x^k\}$, and every $v \in \mathcal{Q}$, there exists a subsequence $J \subseteq \N$ such that, as $j \to \infty$ along $J$, it holds that: $x^j \to x^*,\ v^j_{ex} \to v,\ \norm{v^j_{ex} - v} < \rho,\ \hat{t}^j \in [a,\, b].$
% \[
%   x^j \to x^*,\qquad v^j_{ex} \to v,\qquad \norm{v^j_{ex} - v} < \rho,\qquad \hat{t}^j \in [a,\, b].
% \]
\end{lemma}

\begin{lemma}[Uniform feasibility radius]\label{lem:step-bound}
Let $x^* \in C$ and $v \in \ri(D_C(x^*)) \cap L$. There exist $\rho > 0$, $\bar{\varepsilon} > 0$, and $\delta > 0$ such that: (i)  $B(v,\, \rho) \cap L \subseteq \ri(D_C(x^*))$; (ii) $\tau(x;\, w) \geq \bar{\varepsilon}$ for every $x \in C$ with $\norm{x - x^*} < \delta$ and every $w \in B(v,\, \rho) \cap L$.
% \begin{enumerate}
% \item[(i)] $B(v,\, \rho) \cap L \subseteq \ri(D_C(x^*))$;
% \item[(ii)] $\tau(x;\, w) \geq \bar{\varepsilon}$ for every $x \in C$ with $\norm{x - x^*} < \delta$ and every $w \in B(v,\, \rho) \cap L$.
% \end{enumerate}
\end{lemma}

\begin{lemma}[Continuity and extension]\label{lem:continuity}
The directional derivative map $d \mapsto h'(x^*;\, d)$ is continuous and positively homogeneous of degree one. Moreover, if $h'(x^*;\, v) \geq 0$ for all $v \in \ri(D_C(x^*))$, then $h'(x^*;\, d) \geq 0$ for all $d \in \cl(D_C(x^*))$.
\end{lemma}

\begin{theorem}[Asymptotic d-stationarity]\label{thm:random}
Let  \Cref{ass:A1} and \Cref{ass:A2} hold, and let $\{x^k\}$ be generated by Algorithm~\ref{alg:random} with any oracle $\mathcal{O}_C$ producing feasible iterates. Then:
\begin{enumerate}
\item[(i)] $\sum_{k=0}^{\infty}(t^k_{ex})^2 < \infty$, and in particular $t^k_{ex} \to 0$.
\item[(ii)] Almost surely, every accumulation point of $\{x^k\}$ is d-stationary.
\end{enumerate}
\end{theorem}

\begin{proof}[Proof of \Cref{thm:random}]
\textit{(i) Step-size decay.} Since $t = 0$ is always available in the two-point comparison and $t^k_{ex}$ minimizes over $\{0, \hat{t}^k\}$:
  $h\!\left(x^k + t^k_{ex}\,v^k_{ex}\right) + \tfrac{\gamma}{2}(t^k_{ex})^2 \leq h(x^k).$
  
By greedy selection, $h(x^{k+1}) \leq h(y^{k+1})$, so $h(x^k) - h(x^{k+1}) \geq \frac{\gamma}{2}(t^k_{ex})^2$. Summing gives $\sum_k (t^k_{ex})^2 \leq \frac{2}{\gamma}(h(x^0) - \inf_C h) < \infty$, hence $t^k_{ex} \to 0$.

\textit{(ii) D-stationarity.} Let $\Omega_{BC}$ be the event from Lemma~\ref{lem:BC-directions}, which guarantees simultaneous subsequences for all directions in $\mathcal{Q}$. Fix $\omega \in \Omega_{BC}$ and work pathwise. Let $x^*$ be an accumulation point of $\{x^k(\omega)\}$, and enumerate $\mathcal{Q}\cap \ri(D_C(x^*))$ as $\{v_i\}_{i=0}^{\infty}$. Since $\ri(D_C(x^*))\cap S^{n-1}$ is relatively open in $S^{n-1}\cap L$ and $\mathcal{Q}$ is dense in $S^{n-1}\cap L$, the set $\{v_i\}$ is dense in $\ri(D_C(x^*))\cap S^{n-1}$.

Fix $i \in \N$ and suppose for contradiction $h'(x^*;\, v_i) = -\varepsilon < 0$ for some $\varepsilon > 0$.



By Lemma~\ref{lem:step-bound} applied with $v = v_i$, there exist $\rho_i > 0$, $\bar{\varepsilon}_i > 0$, and $\delta_i > 0$ such that $B(v_i,\, \rho_i) \cap L \subseteq \ri(D_C(x^*))$, and $\tau(x;\, w) \geq \bar{\varepsilon}_i$ for every $x \in C$ with $\norm{x - x^*} < \delta_i$ and every $w \in B(v_i,\, \rho_i) \cap L$.



By local Lipschitz continuity of $h$ at $x^*$, there exist $R_h>0$ and $L_h>0$ such that $h$ is $L_h$-Lipschitz on $B(x^*,R_h)$.

Next choose the comparison scale. By definition of the directional derivative with $h'(x^*;\, v_i) = -\varepsilon$, there exists $\tau_0 > 0$ such that $h(x^* + t v_i) - h(x^*) \leq -\tfrac{\varepsilon}{2}\,t$ for every $t \in (0,\, \tau_0]$. Set
\[
  t^* := \min\!\left\{\bar{\varepsilon}_i,\;\tau_0,\;\frac{\varepsilon}{4\gamma},\;r,\;\frac{R_h}{4}\right\} > 0.
\]



Now lock direction and step-size sampling simultaneously. Since $\omega \in \Omega_{BC}$ and the parameters are now fixed, Lemma~\ref{lem:BC-directions} gives an infinite subsequence $\mathcal{J}_i \subseteq \N$ along which $x^j \to x^*,  v^j_{ex} \to v_i,  \norm{v^j_{ex} - v_i} < \rho_i,  \hat{t}^j \in [t^*/2,\, t^*].$

For $j \in \mathcal{J}_i$ sufficiently large, $\norm{x^j - x^*} < \delta_i$; the feasibility bound above then gives $\tau(x^j;\, v^j_{ex}) \geq \bar{\varepsilon}_i \geq t^*$, so $x^j + t v^j_{ex} \in C$ for every $t \in [0,\, t^*]$.



The key point is that the same strict decrease estimate holds uniformly for every $t \in [t^*/2,\, t^*]$. 
Since $x^j \to x^*$ along $\mathcal{J}_i$, for all sufficiently large $j \in \mathcal{J}_i$ we have $\norm{x^j-x^*}<R_h/4$. For such $j$ and any $t\in[0,t^*]$, $\norm{x^*+t v_i-x^*}=t\le t^*\le R_h/4<R_h,$
% \[
% \norm{x^*+t v_i-x^*}=t\le t^*\le R_h/4<R_h,
% \]
and, using $\norm{v_{ex}^j}=1$,
\[
\norm{x^j+t v_{ex}^j-x^*}\le \norm{x^j-x^*}+t\,\norm{v_{ex}^j}
<R_h/4+t^*\le R_h/2<R_h.
\]
Hence both points $x^j+t v_{ex}^j$ and $x^*+t v_i$ lie in $B(x^*,R_h)$, so the $L_h$-Lipschitz bound is valid. Therefore, for any $t \in [t^*/2,\, t^*]$ and all sufficiently large $j \in \mathcal{J}_i$, decompose
\begin{align*}
  h(x^j + t v^j_{ex}) - h(x^j) + \tfrac{\gamma}{2}t^2
  &= \bigl[h(x^j + t v^j_{ex}) - h(x^* + t v_i)\bigr]
   + \bigl[h(x^* + t v_i) - h(x^*)\bigr] \\
  &\quad + \bigl[h(x^*) - h(x^j)\bigr] + \tfrac{\gamma}{2}t^2 \\
  &\leq L_h\bigl(\norm{x^j - x^*} + t\,\norm{v^j_{ex} - v_i}\bigr)
     - \tfrac{\varepsilon}{2}t
     + L_h\norm{x^j - x^*}
     + \tfrac{\gamma}{2}t^2.
\end{align*}
Since $t \in [t^*/2,\, t^*]$, $\tfrac{\varepsilon}{2}t \geq \tfrac{\varepsilon t^*}{4}$; and $t^* \leq \varepsilon/(4\gamma)$ gives $\tfrac{\gamma}{2}t^2 \leq \tfrac{\gamma}{2}(t^*)^2 \leq \tfrac{\varepsilon t^*}{8}$. 
Since $x^j \to x^*$ and $v^j_{ex} \to v_i$ along $\mathcal{J}_i$, and since $[t^*/2,\, t^*]$ is compact and $x^j \to x^*$, $v^j_{ex} \to v_i$, we may pass to a tail subsequence (still denoted $\mathcal{J}_i$) on which, uniformly in $t \in [t^*/2,\, t^*]$,
\[
  L_h\bigl(\norm{x^j - x^*} + t\,\norm{v^j_{ex} - v_i}\bigr) + L_h\norm{x^j - x^*}
  \;\leq\; 2 L_h\norm{x^j - x^*} + L_h\, t^*\norm{v^j_{ex} - v_i}
  \;\leq\; \tfrac{\varepsilon t^*}{16}.
\]
Therefore, for every $t \in [t^*/2,\, t^*]$ we have that
\[
  h(x^j + t v^j_{ex}) - h(x^j) + \tfrac{\gamma}{2}t^2
  \;\leq\; \tfrac{\varepsilon t^*}{16} - \tfrac{\varepsilon t^*}{4} + \tfrac{\varepsilon t^*}{8}
  \;=\; -\tfrac{\varepsilon t^*}{16} \;<\; 0.
\]
Applied at $t = \hat{t}^j \in [t^*/2,\, t^*]$, together with the feasibility established above, this strict negativity forces the two-point comparison to select $t^j_{ex} = \hat{t}^j \geq t^*/2$ for every $j \in \mathcal{J}_i$.


This yields the contradiction: for infinitely many $j \in \mathcal{J}_i$, $t^j_{ex} \geq t^*/2$, so $\sum_k (t^k_{ex})^2 \geq \sum_{j \in \mathcal{J}_i}(t^*/2)^2 = +\infty$, contradicting part~(i). Hence $h'(x^*;\, v_i) \geq 0$.


Since $\{v_i\}$ is dense in $\ri(D_C(x^*)) \cap S^{n-1}$, Lemma~\ref{lem:continuity} (continuity and positive homogeneity of $h'(x^*;\, \cdot)$, together with the relative-interior-to-closure extension) yields $h'(x^*;\, d) \geq 0$ for every $d \in \cl(D_C(x^*))$. Hence $x^*$ is d-stationary. As $x^*$ was an arbitrary accumulation point on this fixed $\omega \in \Omega_{BC}$, every accumulation point is d-stationary on $\Omega_{BC}$, proving the almost-sure claim.
\end{proof}

\end{document}
